%%%%%%%%%%%%%%%%%%%%%%% Description du fichier %%%%%%%%%%%%%%%%%%%%%%%
% résumé: Préambule d'un template de prise de notes et de rapports.  %
% auteur: Valérien BRAYE, basé sur le travail de SeniorMars (github) %
% dernière modification: 06/10/2023                                  %
%%%%%%%%%%%%%%%%%%%%%%%%%%%%%%%%%%%%%%%%%%%%%%%%%%%%%%%%%%%%%%%%%%%%%%

%%%%%%%%%%%%%%%%%%%%%%%%%%%%%%%%%%%%%%%%%%%%%%%%%%%%%%%%%%%%%%%%%%%%%%%%%%%%%%%
%                               Packages de base                              %
%%%%%%%%%%%%%%%%%%%%%%%%%%%%%%%%%%%%%%%%%%%%%%%%%%%%%%%%%%%%%%%%%%%%%%%%%%%%%%%

% Réduit la taille des marges
\usepackage[a4paper, margin=2cm]{geometry}
% Ajoute plusieurs couleurs.
\usepackage[usenames,dvipsnames,pdftex]{xcolor}
% Permet de change le style des liens.
\usepackage{hyperref}
% Permet d'importer des images et des graphiques.
\usepackage{graphicx}
\graphicspath{{./figures}}
% Permet d'utiliser de nouvelles figures dans des environnements flottants.
\usepackage{float}
% Permet de créer plusieurs colonnes.
\usepackage{multicol}
% Meilleure syntaxe pour les maths.
\usepackage{amsmath,amsfonts,mathtools,amsthm,amssymb}
% Permet de barrer du text.
\usepackage{cancel}
% Permet de modifier le titre des figures.
\usepackage{caption}
% Permet d'importer des pdf directement dans le code LaTeX.
\usepackage{pdfpages}
% Permet d'écrire des algorithmes.
\usepackage[ruled,vlined,linesnumbered]{algorithm2e}
% Change le symbole CQFD (carré blanc en carré noir).
\usepackage{tikzsymbols}
% % Traduction en Français des mots générés (date, table des matières...).
% \usepackage[french]{babel}
% \usepackage[T1]{fontenc}
% Permet l'insertion de texte temporaire
\usepackage{lipsum}
% enumi customs
\usepackage[shortlabels]{enumitem}


% Modifie le symbole CQFD.
\renewcommand\qedsymbol{$\blacksquare$}

\def\class{article}


%%%%%%%%%%%%%%%%%%%%%%%%%%%%%%%%%%%%%%%%%%%%%%%%%%%%%%%%%%%%%%%%%%%%%%%%%%%%%%%
%                              Paramètres de base                             %
%%%%%%%%%%%%%%%%%%%%%%%%%%%%%%%%%%%%%%%%%%%%%%%%%%%%%%%%%%%%%%%%%%%%%%%%%%%%%%%

%%%%%%%%%%%%%%
%  Symboles  %
%%%%%%%%%%%%%%

\let\implies\Rightarrow  % ⟹
\let\impliedby\Leftarrow % ⟸
\let\iff\Leftrightarrow  % ⟺
\let\epsilon\varepsilon  % ε

\AtBeginDocument{\def\labelitemi{$\bullet$}} % change les tirets en points dans les listes
\AtBeginDocument{\def\labelitemii{$\circ$}}

\newcommand*{\lipsumimage}[1][Lorem Ipsum]{
  \begin{figure}[H]
    \centering
    \includegraphics[width=.75\textwidth]{lipsumimage.png}
    \caption{#1}
\end{figure}
}

%%%%%%%%%%%%
%  Tables  %
%%%%%%%%%%%%

\setlength{\tabcolsep}{5pt}
\renewcommand\arraystretch{1.5}

%%%%%%%%%%%%%%
%  SI Unitx  %
%%%%%%%%%%%%%%

\usepackage{siunitx}
\sisetup{locale = FR}

%%%%%%%%%%
%  TikZ  %
%%%%%%%%%%

\usepackage[framemethod=TikZ]{mdframed}
\usepackage{tikz}
\usepackage{tikz-cd}
\usepackage{tikzsymbols}

\usetikzlibrary{intersections, angles, quotes, calc, positioning}
\usetikzlibrary{arrows.meta}

\tikzset{
  force/.style={thick, {Circle[length=2pt]}-stealth, shorten <=-1pt}
}

%%%%%%%%%%%%%%%
%  PGF Plots  %
%%%%%%%%%%%%%%%

\usepackage{pgfplots}
\pgfplotsset{compat=1.13}

%%%%%%%%%%%%%%%%%%%%%%%%%%%%%%%%%%
%  Centrage de la page de titre  %
%%%%%%%%%%%%%%%%%%%%%%%%%%%%%%%%%%

\usepackage{titling}
\renewcommand\maketitlehooka{\null\mbox{}\vfill}
\renewcommand\maketitlehookd{\vfill\null}

%%%%%%%%%%%%%%%%%%%%%%%%%%%%%%%%%%%%%%%%%%%%%%%%%%
%  Crée un arrière plan gris au milieu d'un PDF  %
%%%%%%%%%%%%%%%%%%%%%%%%%%%%%%%%%%%%%%%%%%%%%%%%%%

\usepackage{eso-pic}
\newcommand\definegraybackground{
  \definecolor{reallylightgray}{HTML}{FAFAFA}
  \AddToShipoutPicture{
    \ifthenelse{\isodd{\thepage}}{
      \AtPageLowerLeft{
        \put(\LenToUnit{\dimexpr\paperwidth-222pt},0){
          \color{reallylightgray}\rule{222pt}{297mm}
        }
      }
    }
    {
      \AtPageLowerLeft{
        \color{reallylightgray}\rule{222pt}{297mm}
      }
    }
  }
}

%%%%%%%%%%%%%%%%%%%%%%%%%%%%%%%%%%
%  Modifie la couleur des liens  %
%%%%%%%%%%%%%%%%%%%%%%%%%%%%%%%%%%

\hypersetup{
  % Active la mise en valeur des liens.
  colorlinks,
  % Change la couleur des liens en bleu.
  linkcolor=blue,
  % Change la couleur des citations en noir.
  citecolor={black},
  % Change la couleur des URL de bleu avec un peu de noir
  urlcolor={blue!80!black}
}

%%%%%%%%%%%%%%%%%%
% Fix WrapFigure %
%%%%%%%%%%%%%%%%%%

\newcommand{\wrapfill}{\par\ifnum\value{WF@wrappedlines}>0
    \parskip=0pt
    \addtocounter{WF@wrappedlines}{-1}%
    \null\vspace{\arabic{WF@wrappedlines}\baselineskip}%
    \WFclear
\fi}

%%%%%%%%%%%%%%%%%%
% Multi Colonnes %
%%%%%%%%%%%%%%%%%%

\let\multicolmulticols\multicols
\let\endmulticolmulticols\endmulticols

\RenewDocumentEnvironment{multicols}{mO{}}
{%
  \ifnum#1=1
    #2%
  \else % Plus qu'une colonne
    \multicolmulticols{#1}[#2]
  \fi
}
{%
  \ifnum#1=1
\else % Plus qu'une colonne
  \endmulticolmulticols
\fi
}

\newlength{\thickarrayrulewidth}
\setlength{\thickarrayrulewidth}{5\arrayrulewidth}


%%%%%%%%%%%%%%%%%%%%%%%%%%%%%%%%%%%%%%%%%%%%%%%%%%%%%%%%%%%%%%%%%%%%%%%%%%%
%                     Commandes spécifiques à l'école                     %
%%%%%%%%%%%%%%%%%%%%%%%%%%%%%%%%%%%%%%%%%%%%%%%%%%%%%%%%%%%%%%%%%%%%%%%%%%%

%%%%%%%%%%%%%%%%%%%%%%%%%%%%%%%%%%%%
%  Initialise le nouveau compteur  %
%%%%%%%%%%%%%%%%%%%%%%%%%%%%%%%%%%%%

\newcounter{lecturecounter}

%%%%%%%%%%%%%%%%%%%%%%%%%
%  NewCommandes utiles  %
%%%%%%%%%%%%%%%%%%%%%%%%%

\makeatletter

\newcommand\resetcounters{
  % Réinitialise les compteurs de sous section, sous sous section et la
  % définition de tous les environnements personnalisés.
  \setcounter{subsection}{0}
  \setcounter{subsubsection}{0}
  \setcounter{paragraph}{0}
  \setcounter{subparagraph}{0}
  \setcounter{theorem}{0}
  \setcounter{claim}{0}
  \setcounter{corollary}{0}
  \setcounter{lemma}{0}
  \setcounter{exercise}{0}

  \@ifclasswith\class{nocolor}{
    \setcounter{definition}{0}
  }{}
}

%%%%%%%%%%%%%%%%%%%%%%%%%%%
%  Commandes de chapitre  %
%%%%%%%%%%%%%%%%%%%%%%%%%%%

\usepackage{xifthen}

% EXEMPLE:
% 1. \lesson{Date (Heure)}{Titre du chapitre}
% 2. \lesson[4]{Date (Heure)}{Titre du chapitre}
% 3. \lesson{Date (Titre du chapitre)}{}
% 4. \lesson[4]{Date (Heure)}{}
% Paramètres:
% 1. (Optionnel) Numéro du chapitre.
% 2. Date et heure du chapitre.
% 3. Titre du chapitre.
\def\@lesson{}
\newcommand\lesson[3][\arabic{lecturecounter}]{
  % Incrémente le lecturecounter.
  \addtocounter{lecturecounter}{1}

  % Met le numéro de section au numéro de chapitre.
  \setcounter{section}{#1}
  \renewcommand\thesubsection{#1.\arabic{subsection}}

  % Réinitialise les compteurs.
  \resetcounters

  % Vérifie si l'utilisateur a fourni le titre du chapitre.
  \ifthenelse{\isempty{#3}}{
    \def\@lesson{Chapitre \arabic{lecturecounter}}
  }{
    \def\@lesson{Chapitre \arabic{lecturecounter}: #3}
  }

  % Affiche l'information sous la forme suivante:
  %                                                                Date (Heure)
  % ---------------------------------------------------------------------------
  % Chapitre 1: Titre du chapitre
  \hfill\small{#2}
  \hrule
  \vspace*{-0.3cm}
  \section*{\@lesson}
  \addcontentsline{toc}{section}{\@lesson}
}

%%%%%%%%%%%%%%%%%%%%%%%%
%  Import des figures  %
%%%%%%%%%%%%%%%%%%%%%%%%

\usepackage{import}
\pdfminorversion=7

% EXEMPLE:
% 1. \incfig{limit-graph}
% 2. \incfig[0.4]{limit-graph}
% Paramètres:
% 1. Le nom de la figure. Elle doit être dans figures/NOM.tex_pdf.
% 2. (Optionnel) La largeur de la figure. Exemple: 0.5, 0.35.
\newcommand\incfig[2][1]{%
  \def\svgwidth{#1\columnwidth}
  \import{./figures/}{#2.pdf_tex}
}

\begingroup\expandafter\expandafter\expandafter\endgroup
\expandafter\ifx\csname pdfsuppresswarningpagegroup\endcsname\relax
\else
  \pdfsuppresswarningpagegroup=1\relax
\fi

%%%%%%%%%%%%%%%%%
% Fancy Headers %
%%%%%%%%%%%%%%%%%

\usepackage{fancyhdr}

% Force une nouvelle page.
\newcommand\forcenewpage{\clearpage\mbox{~}\clearpage\newpage}

% Cette commande simplifie la gestion des entêtes et pieds de page.
\newcommand\createintro{
  % Utilise les chiffres romains (ie. i, v, vi, x, ...).
  \pagenumbering{roman}

  % Affiche le style de la page.
  \maketitle
  % Rend le titre du pagestyle vide, sans fancy headers / footers.
  \thispagestyle{empty}
  % Saute une page
  \newpage


  % Rend le pagestyle fancy pour la page intro.tex.
  \pagestyle{fancy}

  % Retire la ligne de l'entête
  \renewcommand\headrulewidth{0pt}

  % Retire tout ce qui est lié à l'entête.
  \fancyhead{}

  % Ajoute ce qui est lié à l'entête au centre.
  \fancyfoot[C]{
    \vspace{-2cm} % remonte le pied de page (sinon marge trop petite)
    \hrule
    \vspace{0.1cm}
    \@author \\
    \term \space \academicyear \\
    Dernière mise à jour: \@date \\
    \faculty
  }

  \newpage

  % Change la numérotation des pages en chiffres arabes.
  \pagenumbering{arabic}
  % Règle le numéro de page à 1.
  \setcounter{page}{1}

  % Ajoute la ligne de l'entête précédemment enlevée.
  \renewcommand\headrulewidth{0.4pt}
  % Ajoute le titre du chapitre en haut à droite.
  \fancyhead[R]{\@lesson}
  % Ajoute le titre de l'auteur en haut à gauche.
  \fancyhead[L]{\@author}
  % Ajoute le numéro de page en bas au centre.
  \fancyfoot[C]{\thepage}
}

\makeatother


%%%%%%%%%%%%%%%%%%%%%%%%%%%%%%%%%%%%%%%%%%%%%%%%%%%%%%%%%%%%%%%%%%%%%%%%%%%%%%%%%%
%                            Commandes personnalisées                            %
%%%%%%%%%%%%%%%%%%%%%%%%%%%%%%%%%%%%%%%%%%%%%%%%%%%%%%%%%%%%%%%%%%%%%%%%%%%%%%%%%%

%%%%%%%%%%%%%%
%  Encercle  %
%%%%%%%%%%%%%%

\newcommand*\circled[1]{\tikz[baseline=(char.base)]{
  \node[shape=circle,draw,inner sep=1pt] (char) {#1};}
}

%%%%%%%%%%%%%%%%%%%%
%  Commandes TODO  %
%%%%%%%%%%%%%%%%%%%%

\usepackage{xargs}
\usepackage[colorinlistoftodos]{todonotes}

\makeatletter

\@ifclasswith\class{working}{
  \newcommandx\unsure[2][1=]{\todo[linecolor=red,backgroundcolor=red!25,bordercolor=red,#1]{#2}}
  \newcommandx\change[2][1=]{\todo[linecolor=blue,backgroundcolor=blue!25,bordercolor=blue,#1]{#2}}
  \newcommandx\info[2][1=]{\todo[linecolor=OliveGreen,backgroundcolor=OliveGreen!25,bordercolor=OliveGreen,#1]{#2}}
  \newcommandx\improvement[2][1=]{\todo[linecolor=Plum,backgroundcolor=Plum!25,bordercolor=Plum,#1]{#2}}

  \newcommand\listnotes{
    \newpage
    \listoftodos[Notes]
  }
}{
  \newcommandx\unsure[2][1=]{}
  \newcommandx\change[2][1=]{}
  \newcommandx\info[2][1=]{}
  \newcommandx\improvement[2][1=]{}

  \newcommand\listnotes{}
}

\makeatother

%%%%%%%%%%%%%
%  Correct  %
%%%%%%%%%%%%%

% EXEMPLE:
% 1. \correct{INCORRECT}{CORRECT}
% Paramètres:
% 1. L'affirmation incorrecte.
% 2. L'affirmation correcte.
\definecolor{correct}{HTML}{009900}
\newcommand\correct[2]{{\color{red}{#1 }}\ensuremath{\to}{\color{correct}{ #2}}}


%%%%%%%%%%%%%%%%%%%%%%%%%%%%%%%%%%%%%%%%%%%%%%%%%%%%%%%%%%%%%%%%%%%%%%%%%%%%%%%
%                                Environnements                               %
%%%%%%%%%%%%%%%%%%%%%%%%%%%%%%%%%%%%%%%%%%%%%%%%%%%%%%%%%%%%%%%%%%%%%%%%%%%%%%%

\usepackage{varwidth}
\usepackage{thmtools}
\usepackage[most,many,breakable]{tcolorbox}

\tcbuselibrary{theorems,skins,hooks}
\usetikzlibrary{arrows,calc,shadows.blur}

%%%%%%%%%%%%%%%%%%%%%%%%%%%%%%
%  Définitions des couleurs  %
%%%%%%%%%%%%%%%%%%%%%%%%%%%%%%

\definecolor{myblue}{RGB}{45, 111, 177}
\definecolor{mygreen}{RGB}{56, 140, 70}
\definecolor{myred}{RGB}{199, 68, 64}
\definecolor{mypurple}{RGB}{197, 92, 212}

\definecolor{definition}{HTML}{228b22}
\definecolor{theorem}{HTML}{00007B}
\definecolor{example}{HTML}{2A7F7F}
\definecolor{definition}{HTML}{228b22}
\definecolor{prop}{HTML}{191971}
\definecolor{lemma}{HTML}{983b0f}
\definecolor{exercise}{HTML}{88D6D1}

\colorlet{definition}{mygreen!85!black}
\colorlet{claim}{mygreen!85!black}
\colorlet{corollary}{mypurple!85!black}
\colorlet{proof}{theorem}

%%%%%%%%%%%%%%%%%%%%%%%%%%%%%%%%%%%%%%%%%%%%%%%%%%%%%%%%%%%%%%%
%  Crée des environnements de style basés sur les paramètres  %
%%%%%%%%%%%%%%%%%%%%%%%%%%%%%%%%%%%%%%%%%%%%%%%%%%%%%%%%%%%%%%%

\mdfsetup{skipabove=1em,skipbelow=0em}

%%%%%%%%%%%%%%%%%%%%%%
%  Commandes utiles  %
%%%%%%%%%%%%%%%%%%%%%%

% EXEMPLE:
% 1. \createnewtheoremstyle{thmdefinitionbox}{}{}
% 2. \createnewtheoremstyle{thmtheorembox}{}{}
% 3. \createnewtheoremstyle{thmproofbox}{qed=\qedsymbol}{
%       rightline=false, topline=false, bottomline=false
%    }
% Paramètres:
% 1. Nom du théorème.
% 2. N'importe quels paramètres supplémentaires pour declaretheoremstyle.
% 3. N'importe quels paramètres supplémentaires pour mdframed.
\newcommand\createnewtheoremstyle[3]{
  \declaretheoremstyle[
  headfont=\bfseries\sffamily, bodyfont=\normalfont, #2,
  mdframed={
    #3,
  },
  ]{#1}
}

% EXEMPLE:
% 1. \createnewcoloredtheoremstyle{thmdefinitionbox}{definition}{}{}
% 2. \createnewcoloredtheoremstyle{thmexamplebox}{example}{}{
%       rightline=true, leftline=true, topline=true, bottomline=true
%     }
% 3. \createnewcoloredtheoremstyle{thmproofbox}{proof}{qed=\qedsymbol}{backgroundcolor=white}
% Paramètres:
% 1. Nom du théorème.
% 2. Color of theorem.
% 2. N'importe quels paramètres supplémentaires pour declaretheoremstyle.
% 3. N'importe quels paramètres supplémentaires pour mdframed.
\newcommand\createnewcoloredtheoremstyle[4]{
  \declaretheoremstyle[
  headfont=\bfseries\sffamily\color{#2}, bodyfont=\normalfont, #3,
  mdframed={
    linewidth=2pt,
    rightline=false, leftline=true, topline=false, bottomline=false,
    linecolor=#2, backgroundcolor=#2!5, #4,
  },
  ]{#1}
}

%%%%%%%%%%%%%%%%%%%%%%%%%%%%%%%%%%%%%%
%  Crée les environnements de style  %
%%%%%%%%%%%%%%%%%%%%%%%%%%%%%%%%%%%%%%

\makeatletter
\@ifclasswith\class{nocolor}{
  % Environnements sans couleurs.

  \createnewtheoremstyle{thmdefinitionbox}{}{}
  \createnewtheoremstyle{thmtheorembox}{}{}
  \createnewtheoremstyle{thmexamplebox}{}{}
  \createnewtheoremstyle{thmclaimbox}{}{}
  \createnewtheoremstyle{thmcorollarybox}{}{}
  \createnewtheoremstyle{thmpropbox}{}{}
  \createnewtheoremstyle{thmlemmabox}{}{}
  \createnewtheoremstyle{thmexercisebox}{}{}
  \createnewtheoremstyle{thmdefinitionbox}{}{}
  \createnewtheoremstyle{thmquestionbox}{}{}
  \createnewtheoremstyle{thmsolutionbox}{}{}

  \createnewtheoremstyle{thmproofbox}{qed=\qedsymbol}{}
  \createnewtheoremstyle{thmexplanationbox}{}{}
}{
  % Environnements avec couleurs.

  \createnewcoloredtheoremstyle{thmdefinitionbox}{definition}{}{}
  \createnewcoloredtheoremstyle{thmtheorembox}{theorem}{}{}
  \createnewcoloredtheoremstyle{thmexamplebox}{example}{}{
    rightline=true, leftline=true, topline=true, bottomline=true
  }
  \createnewcoloredtheoremstyle{thmclaimbox}{claim}{}{}
  \createnewcoloredtheoremstyle{thmcorollarybox}{corollary}{}{}
  \createnewcoloredtheoremstyle{thmpropbox}{prop}{}{}
  \createnewcoloredtheoremstyle{thmlemmabox}{lemma}{}{}
  \createnewcoloredtheoremstyle{thmexercisebox}{exercise}{}{}

  \createnewcoloredtheoremstyle{thmproofbox}{proof}{qed=\qedsymbol}{backgroundcolor=white}
  \createnewcoloredtheoremstyle{thmexplanationbox}{example}{qed=\qedsymbol}{backgroundcolor=white}
}
\makeatother

%%%%%%%%%%%%%%%%%%%%%%%%%%%%%
%  Crée les environnements  %
%%%%%%%%%%%%%%%%%%%%%%%%%%%%%

\declaretheorem[numberwithin=section, style=thmtheorembox,     name=Théorème]{theorem}
\declaretheorem[numbered=no,          style=thmexamplebox,     name=Bonus]{bonus}
\declaretheorem[numberwithin=section, style=thmclaimbox,       name=Affirmation]{claim}
\declaretheorem[numberwithin=section, style=thmcorollarybox,   name=Corollaire]{corollary}
\declaretheorem[numberwithin=section, style=thmpropbox,        name=Proposition]{prop}
\declaretheorem[numberwithin=section, style=thmlemmabox,       name=Lemme]{lemma}
\declaretheorem[numberwithin=section, style=thmexercisebox,    name=Exercice]{exercise}
\declaretheorem[numbered=no,          style=thmproofbox,       name=Preuve]{replacementproof}
\declaretheorem[numbered=no,          style=thmexplanationbox, name=Preuve]{expl}

\makeatletter
\@ifclasswith\class{nocolor}{
  % Environnements sans couleurs.

  \newtheorem*{note}{Note}

  \declaretheorem[numbered=no, style=thmdefinitionbox, name=Definition]{definition}
  \declaretheorem[numberwithin=section, style=thmquestionbox,   name=Question]{question}
  \declaretheorem[numberwithin=section, style=thmsolutionbox,   name=Solution]{solution}
}{
  % Environnements avec couleurs.

  \newtcbtheorem{Definition}{Definition}{
    enhanced,
    before skip=2mm,
    after skip=2mm,
    colback=red!5,
    colframe=red!80!black,
    colbacktitle=red!75!black,
    boxrule=0.5mm,
    attach boxed title to top left={
      xshift=1cm,
      yshift*=1mm-\tcboxedtitleheight
    },
    varwidth boxed title*=-3cm,
    boxed title style={
      interior engine=empty,
      frame code={
        \path[fill=tcbcolback]
        ([yshift=-1mm,xshift=-1mm]frame.north west)
        arc[start angle=0,end angle=180,radius=1mm]
        ([yshift=-1mm,xshift=1mm]frame.north east)
        arc[start angle=180,end angle=0,radius=1mm];
        \path[left color=tcbcolback!60!black,right color=tcbcolback!60!black,
        middle color=tcbcolback!80!black]
        ([xshift=-2mm]frame.north west) -- ([xshift=2mm]frame.north east)
        [rounded corners=1mm]-- ([xshift=1mm,yshift=-1mm]frame.north east)
        -- (frame.south east) -- (frame.south west)
        -- ([xshift=-1mm,yshift=-1mm]frame.north west)
        [sharp corners]-- cycle;
      },
    },
    fonttitle=\bfseries,
    title={#2},
    #1
  }{def}

  \NewDocumentEnvironment{definition}{O{}O{}}
    {\begin{Definition}{#1}{#2}}{\end{Definition}}

  \newtcolorbox{note}[1][]{%
    enhanced jigsaw,
    colback=gray!20!white,%
    colframe=gray!80!black,
    size=small,
    boxrule=1pt,
    title=\textbf{Note: Source du style des graphiques},
    halign title=flush center,
    coltitle=black,
    breakable,
    drop shadow=black!50!white,
    attach boxed title to top left={xshift=1cm,yshift=-\tcboxedtitleheight/2,yshifttext=-\tcboxedtitleheight/2},
    minipage boxed title=6.75cm,
    boxed title style={%
      colback=white,
      size=fbox,
      boxrule=1pt,
      boxsep=2pt,
      underlay={%
        \coordinate (dotA) at ($(interior.west) + (-0.5pt,0)$);
        \coordinate (dotB) at ($(interior.east) + (.5pt,0)$);
        \begin{scope}
          \clip (interior.north west) rectangle ([xshift=3ex]interior.east);
          \filldraw [white, blur shadow={shadow opacity=60, shadow yshift=-.75ex}, rounded corners=2pt] (interior.north west) rectangle (interior.south east);
        \end{scope}
        \begin{scope}[gray!80!black]
          \fill (dotA) circle (2pt);
          \fill (dotB) circle (2pt);
        \end{scope}
      },
    },
    #1,
  }

  \newtcbtheorem{Question}{Question}{enhanced,
    breakable,
    colback=white,
    colframe=myblue!80!black,
    attach boxed title to top left={yshift*=-\tcboxedtitleheight},
    fonttitle=\bfseries,
    title=\textbf{Question:},
    boxed title size=title,
    boxed title style={
      sharp corners,
      rounded corners=northwest,
      colback=tcbcolframe,
      boxrule=0pt,
    },
    underlay boxed title={
      \path[fill=tcbcolframe] (title.south west)--(title.south east)
      to[out=0, in=180] ([xshift=5mm]title.east)--
      (title.center-|frame.east)
      [rounded corners=\kvtcb@arc] |-
      (frame.north) -| cycle;
    },
    #1
  }{def}

  \NewDocumentEnvironment{question}{O{}O{}}
  {\begin{Question}{#1}{#2}}{\end{Question}}

  \newtcolorbox{Solution}{enhanced,
    breakable,
    colback=white,
    colframe=mygreen!80!black,
    attach boxed title to top left={yshift*=-\tcboxedtitleheight},
    title=\textbf{Solution:},
    boxed title size=title,
    boxed title style={
      sharp corners,
      rounded corners=northwest,
      colback=tcbcolframe,
      boxrule=0pt,
    },
    underlay boxed title={
      \path[fill=tcbcolframe] (title.south west)--(title.south east)
      to[out=0, in=180] ([xshift=5mm]title.east)--
      (title.center-|frame.east)
      [rounded corners=\kvtcb@arc] |-
      (frame.north) -| cycle;
    },
  }

  \NewDocumentEnvironment{solution}{O{}O{}}
  {\vspace{-10pt}\begin{Solution}{#1}{#2}}{\end{Solution}}
}
\makeatother

%%%%%%%%%%%%%%%%%%%%%%%%%%%%%%%%%%%%%%%%
%  Environnements d'édition de preuve  %
%%%%%%%%%%%%%%%%%%%%%%%%%%%%%%%%%%%%%%%%

\renewenvironment{proof}[1][\proofname]{\vspace{-10pt}\begin{replacementproof}}{\end{replacementproof}}
\newenvironment{explanation}[1][\proofname]{\vspace{-10pt}\begin{expl}}{\end{expl}}

\theoremstyle{definition}

\newtheorem*{notation}{Notation}
\newtheorem*{weknowthat}{We know that}
\newtheorem*{problem}{Problem}
\newtheorem*{observe}{We observe that}
\newtheorem*{property}{Property}
\newtheorem*{intuition}{Intuition}
\newtheorem*{remark}{Remark}